% Solutions of Book 4, chapter 9 (functional-programming).
\section*{Chapter \ref{ch:b2:functional-programming} --- Functional Programming}

\begin{solution}{exo:b2:functional-programming:1}
\omoutput{code/bachelor-2/09-functional-programming/outputs/exo_lazy_order.txt}{The order.}
Nothing runs before \code{sum}. Then 1 goes through both \code{inspect}s, becomes 10 and is refused by the \omterm{def:b2:functional-programming:fold}{filter}; 2
becomes 20, passes, and \code{take(1)} has its element and stops asking. Two elements of \code{v} are read, and 3 and 4
never.
\end{solution}

\begin{solution}{exo:b2:functional-programming:2}
\omoutput[2-6,15-16]{code/bachelor-2/09-functional-programming/outputs/exo_fnmut.txt}{The error, shortened.}
The \omterm{def:b2:functional-programming:closure}{closure} adds to \code{late}, so it implements \code{FnMut} only. The parameter becomes \code{mut f: impl FnMut(u32)}. \code{Fn} promises that the \omterm{def:b2:functional-programming:closure}{closure} can be called through a shared reference,
possibly from several places at once; a \omterm{def:b2:functional-programming:closure}{closure} that modifies a captured variable needs exclusive access during
each call, which only \code{FnMut} (called through \code{&mut}) gives. The program then prints \code{3 late}.
\end{solution}

\begin{solution}{exo:b2:functional-programming:3}
\omcode{code/bachelor-2/09-functional-programming/exercises/compose.rs}{2}{9}{Composition}
The \omterm{def:b2:functional-programming:closure}{closure} returned by \code{compose} must own \code{f} and \code{g}, hence \code{move}. \code{compose_all} \omterm{def:b2:functional-programming:fold}{folds}
over the functions from the last to the first, so \code{compose_all(vec![f, g, h])(x)} is $f(g(h(x)))$; the test
checks $((2 + 1)^2) - 3 = 6$.
\end{solution}

\begin{solution}{exo:b2:functional-programming:4}
\omoutput{code/bachelor-2/09-functional-programming/outputs/exo_sizes.txt}{Sizes, then values.}
\code{a} captures nothing: 0 bytes, like a plain function. \code{b} reads \code{rate} and \code{title} and captures
them by reference: two pointers, 16 bytes. \code{c} moves in the array of four \code{u64} (32 bytes) and the
\code{u32}, 36 bytes, padded to 40 for the alignment of the \code{u64}s.
\end{solution}

\begin{solution}{exo:b2:functional-programming:5}
\omoutput{code/bachelor-2/09-functional-programming/outputs/exo_folds.txt}{The four values.}
\code{a} reads the digits in order, 31415; \code{b} in reverse, 51413; \code{c} inserts each element at the front,
reversing the list; \code{d} keeps the minimum and maximum, \code{(1, 5)}. A left \omterm{def:b2:functional-programming:fold}{fold} over the reversed list is a right
\omterm{def:b2:functional-programming:fold}{fold} with the arguments swapped, so \code{b} is the right \omterm{def:b2:functional-programming:fold}{fold} of \code{a}'s function. Only \code{d} gives the same
result from either side, because \code{min} and \code{max} are associative and commutative; the right \omterm{def:b2:functional-programming:fold}{fold} of
\code{c}'s function rebuilds the list in its original order.
\end{solution}

\begin{solution}{exo:b2:functional-programming:6}
\omoutput[2-9]{code/bachelor-2/09-functional-programming/outputs/exo_twice.txt}{The error, shortened.}
\code{count} takes the iterator by value and runs it to the end, so \code{long} no longer exists for the loop; an
iterator is a cursor, not a collection. Fix 1: a \omterm{def:b2:functional-programming:closure}{closure} that builds it, \code{let long = || titles.iter().filter(|t|
t.len() > 13);}, then \code{long().count()} and \code{for t in long()}; the \omterm{def:b2:functional-programming:fold}{filter} runs twice. Fix 2:
\code{let long: Vec<_> = ... .collect();}, then \code{long.len()} and \code{for t in &long}; the \omterm{def:b2:functional-programming:fold}{filter} runs once, at
the cost of a vector.
\end{solution}

\begin{solution}{exo:b2:functional-programming:7}
\omcode{code/bachelor-2/09-functional-programming/exercises/word_freq.rs}{4}{19}{Frequencies by one \omterm{def:b2:functional-programming:fold}{fold}}
The \omterm{def:b2:functional-programming:fold}{fold}'s accumulator is the map itself, \omterm{def:b2:rust-ownership:move}{moved from} step to step. \code{then_with} breaks ties in the count by
the word. The test checks a sentence in which \emph{the} appears four times in three spellings.
\end{solution}

\begin{solution}{exo:b2:functional-programming:8}
For $[\,]$: $\mathrm{foldr}\,(:)\,[\,]\,[\,] = [\,]$. For $x : xs$: $\mathrm{foldr}\,(:)\,[\,]\,(x : xs) = x :
\mathrm{foldr}\,(:)\,[\,]\,xs = x : xs$, by the hypothesis. The left \omterm{def:b2:functional-programming:fold}{fold} of $(\lambda acc\,x.\ x : acc)$ from
$[\,]$ puts each element in front of the previous ones: it reverses the list (\code{c} of
\cref{exo:b2:functional-programming:5}). The crate checks the first law on 500 lists.
\end{solution}

\begin{solution}{exo:b2:functional-programming:9}
\omcode{code/bachelor-2/09-functional-programming/exercises/queue.rs}{32}{50}{Push, and pop with a reversal}
The back list holds the most recent element first, so the front list is its reversal. \code{pop} returns, besides
the element and the new queue, the number of nodes copied by a reversal, which the tests and
\cref{exo:b2:functional-programming:12} use.
\end{solution}

\begin{solution}{exo:b2:functional-programming:10}
\omterm{def:b2:functional-programming:fold}{Filter} fusion, by induction on $xs$. For $[\,]$ both sides are $[\,]$. For $x : xs$: if $q\,x$ is false, the inner
\omterm{def:b2:functional-programming:fold}{filter} drops $x$ and the left side is $\mathrm{filter}\,p\,(\mathrm{filter}\,q\,xs)$; the right side drops $x$ too,
since $q\,x \wedge p\,x$ is false, and the hypothesis closes the case. If $q\,x$ is true, the left side is
$\mathrm{filter}\,p\,(x : \mathrm{filter}\,q\,xs)$, which keeps $x$ exactly when $p\,x$, as the right side does, and
the tails agree by the hypothesis. The second law: if $p\,(f\,x)$, the left side is $f\,x : \mathrm{map}\,f\,
(\mathrm{filter}\,(p \circ f)\,xs)$ and the right side $\mathrm{filter}\,p\,(f\,x : \mathrm{map}\,f\,xs) = f\,x :
\mathrm{filter}\,p\,(\mathrm{map}\,f\,xs)$; otherwise both drop the element; the hypothesis gives the tails. With \omterm{def:b2:functional-programming:pure}{side
effects} the laws fail as programs: the left side of the first calls $p$ only on elements that pass $q$, and the second
law calls $f$ once or twice per element depending on the side.
\end{solution}

\begin{solution}{exo:b2:functional-programming:11}
\omcode{code/bachelor-2/09-functional-programming/exercises/pmap.rs}{23}{39}{Insertion that replaces in the copy}
Each step copies the node it passes through, sharing the untouched child; at the key, the copy takes the new value,
and the old node keeps the old one. The test inserts 1000 pairs over 100 keys and compares every 37th version, key by
key, with the \code{BTreeMap} snapshot of the same step.
\end{solution}

\begin{solution}{exo:b2:functional-programming:12}
Push $n$ elements into an empty queue: all are in the back list, the front is empty; call this version $q$. Now pop
$q$, keep the element, throw away the new queue, and pop $q$ again, $m$ times: each pop finds the front of $q$ empty
and reverses its back list of $n$ nodes, $n$ steps, since $q$ itself never changes. Total $mn$; the test checks two
pops of an old version costing 11 each. The amortised argument charges each element's move to the push that brought
it, and assumes that a move, once paid, is never needed again; with persistence the same old version, with its
unpaid reversal, can be used any number of times.
\end{solution}

\begin{solution}{pb:b2:functional-programming:1}
\textbf{1.} \code{Rc<dyn Fn(&Sheet) -> f64>}. The sheet is a parameter so that the formula reads the version it is
evaluated in: a captured sheet would be one fixed version, and the sheet would contain a formula owning the sheet, a
cycle of \code{Rc} that is never freed.

\textbf{2.} \code{set} clones the cell of every node it copies (\code{n.cell.clone()}). A \code{Box<dyn Fn>} cannot be
cloned; an \code{Rc} clones by incrementing a count, sharing the \omterm{def:b2:functional-programming:closure}{closure}.

\textbf{3.} \code{Fn}: a formula is called any number of times, through shared references, from any version.
\code{FnMut} would need exclusive access to the formula; \code{FnOnce} allows a single call.

\textbf{4.} The column number \code{c}, by value. \code{c} is the loop variable and dies at the end of each turn of
the loop; the \omterm{def:b2:functional-programming:closure}{closure} lives in the sheet, so it must own its copy.

\textbf{5.} Yes: its result depends only on the cells of the sheet it is given. A formula that modified a captured
variable would give results that depend on how often and in which order cells were evaluated, which laziness makes
unpredictable; \code{Fn} forbids it anyway.

\textbf{6.} It calls \code{value} on its own cell, which calls the formula again: an infinite recursion, ended by a
stack overflow. A real sheet detects cycles before evaluating.

\textbf{7.} None: \emph{formulas run: 0}.

\textbf{8.} One formula; it reads nine cells, \code{5, 15, ..., 85}.

\textbf{9.} \code{value} takes \code{&self}, a shared reference, yet must change the counter; \code{Cell} allows a
change through a shared reference (for a \code{Copy} value, by \code{get} and \code{set}). The \code{Rc} shares one
counter among all versions.

\textbf{10.} Version 0 still holds its original cells, whose column total is 27. Version 1000 holds the values written
by the edits; the formula, the same \omterm{def:b2:functional-programming:closure}{closure}, reads those.

\textbf{11.} Store, with each formula node, the value computed for a given version. Since a version never changes, a
value cached for a version stays valid forever for it; it is valid in a later version too if every cell the formula
reads is the same node in both (a pointer comparison), which \omterm{def:b2:functional-programming:persistent}{path copying} makes cheap to check.

\textbf{12.} Building a version costs $O(\log n)$ whatever its formulas; a version that is never displayed never runs
them.

\textbf{13.} The key exists at depth $d$, so the path from the root to it has $d + 1$ nodes, and each is copied once;
the node of the cell is replaced by a copy holding the new cell. $d + 1$ new nodes, as in
\cref{prop:b2:functional-programming:path-copy}, and nothing is added to the tree.

\textbf{14.} Let $L(n)$ be the number of levels for $n$ keys. $L(1) = 1$, and for $n \ge 2$ the larger part has
$\lfloor n/2 \rfloor$ keys, so $L(n) = 1 + L(\lfloor n/2 \rfloor)$. By induction, $L(n) \le 1 + \lfloor \log_2 \lfloor
n/2 \rfloor \rfloor + 1 = \lfloor \log_2 n \rfloor + 1$, since $\lfloor \log_2 \lfloor n/2 \rfloor \rfloor =
\lfloor \log_2 n \rfloor - 1$ for $n \ge 2$.

\textbf{15.} A \code{Num} copies 16 bytes (the cell); a \code{Formula} copies a \omterm{def:b2:rust-traits-and-generics:trait-object}{fat pointer} and increments the
\code{Rc} count: no \omterm{def:b2:functional-programming:closure}{closure} is copied.

\textbf{16.} \code{seen.insert} returns \code{false} for a node already counted, and the walk stops there, so each
distinct node is entered once and each version costs one more step for its root.

\textbf{17.} Exactly the nodes none of whose descendants is an edited cell: every number cell is edited (question
18), so a node is shared only if its whole subtree holds formula cells, keys 90 to 99. For this tree there are 7 such
nodes; all others were copied at least once between the two versions.

\textbf{18.} 37 and 90 have no common divisor, so $i \mapsto 37i \bmod 90$ runs through all 90 residues every 90
edits. $1000 = 11 \cdot 90 + 10$: 10 cells are edited 12 times, the other 80 cells 11 times.

\textbf{19.} The initial tree has $n$ nodes. By question 13, an edit at depth $d$ creates $d + 1$ nodes, and by
question 14, $d + 1 \le \lfloor \log_2 n \rfloor + 1 \le \lceil \log_2 n \rceil + 1$ (edits never change the shape).
So $k$ edits create at most $k(\lceil \log_2 n \rceil + 1)$ nodes, and all versions together hold at most $n + k(\lceil
\log_2 n \rceil + 1)$ distinct nodes.

\textbf{20.} 5891 against 8100. The bound charges every edit for the deepest path, 8 levels by the formula, while
the tree has 7, and only 37 of its 100 nodes are on the last level; the average edit creates about
5.8 nodes.

\textbf{21.} $100100 / 5891 \approx 17$. Full copies grow as $kn$, the \omterm{def:b2:functional-programming:persistent}{persistent} versions as $k \log_2 n$: the ratio
grows like $n / \log_2 n$.

\textbf{22.} $5891 \times 40 = 235{,}640$ bytes, about 230\,KiB. The 40 bytes are the \code{Node}; each
\code{Rc} allocation also holds a strong and a weak count, 16 more bytes, so about 330\,KB in all, against about
5.6\,MB for full copies of the nodes at the same size.

\textbf{23.} \code{versions.push(edited)}: the vector owns every version, so no root's count reaches zero.

\textbf{24.} $k$ versions of an $n$-cell sheet in at most $n + k(\lceil \log_2 n \rceil + 1)$ nodes: for 100 cells
and 1000 edits, 5891 distinct nodes, under the bound of 8100, against 100100 for full copies.
\end{solution}
